\section{System Architecture}
\label{sec:arch}

Three retrieval arms answer the same question over the same stores.
\Cref{fig:arch} shows all three.

\subsection{Text-to-SQL}

The model is shown the complete relational schema, including the
\texttt{subsidiary} and \texttt{reporting\_owner} tables, and writes a SQL query;
the query is executed; the resulting rows are given back to the model, which
writes the answer. Two LLM calls.

This is deliberately a strong baseline rather than a straw man. Because the schema
is complete, every multi-hop answer in the benchmark \emph{is} reachable by joins.
What the model must do unaided is choose the join path, decode a multi-valued
relationship string, and recognize that traversing an ownership chain of unknown
depth requires a recursive CTE. The schema documentation also states the
surname-first convention used for person names --- an asymmetry in an earlier
iteration, where only the graph documentation mentioned it, is recorded among the
harness defects in \Cref{sec:repro}.

\subsection{Vector RAG}

The question is embedded and the $k{=}8$ nearest chunks are retrieved by cosine
similarity from an HNSW index \citep{malkov2018hnsw,pgvector2024}; the chunks are
given to the model, which answers. One LLM call, and no query language, therefore
no repair path.

The encoder is \texttt{all-MiniLM-L6-v2}, a six-layer 22.7M-parameter model
distilled in the MiniLM family \citep{wang2020minilm} and trained as a sentence
encoder \citep{reimers2019sbert}. It emits 384 dimensions with mean pooling and
$L^2$ normalization, so cosine similarity reduces to a dot product.

One configuration detail is load-bearing. The encoder accepts 256 word-pieces,
roughly a thousand characters, while Item 1A averages 116{,}779 characters. The
section therefore cannot be embedded whole: text is chunked at 1{,}400 characters
with 200 characters of overlap, so a single risk-factor section becomes on the
order of a hundred vectors, and 5{,}210 sections become 508{,}714 chunks. We
embed only Item 1A. That is not a handicap imposed on the baseline but the
opposite --- it is the only section any benchmark question asks about, so
restricting the index concentrates the candidate pool and makes the vector arm
\emph{stronger} than it would be with all seven Items embedded.

\subsection{GraphRAG}

Two LLM calls, as with text-to-SQL, but the work between them is split
differently.

\begin{enumerate}
  \item \textbf{Identify.} The model writes a Cypher query whose job is to
        determine \emph{which entities} the question concerns --- typically
        resolving a person to a CIK and finding the companies where they hold a
        role, with their relevant filings.
  \item \textbf{Expand (deterministic).} The system, not the model, expands each
        identified entity: for every company returned, fetch its subsidiaries and
        its risk-factor filings.
  \item \textbf{Fetch text by key.} The accession numbers from the traversal are
        used to select section text by primary key, and windows of $\pm 900$
        characters around matched terms are extracted, at most three per filing.
        This is a lookup, not a search: the graph has already decided which
        documents are in scope, so there is nothing to rank and nothing to miss.
  \item \textbf{Answer.} The model receives the graph facts and the exact
        passages, and writes the answer with citations.
\end{enumerate}

Stage 2 being system-authored rather than LLM-authored is the design's central
claim, and it is a response to a measured failure rather than an aesthetic
preference. An LLM writing a single query over a fan-out join will attach a
\texttt{LIMIT}, and one entity's hundreds of child rows then consume the entire
budget --- after which the model reports, fluently and wrongly, that the other
entities have no subsidiaries or do not exist. Splitting identification from
expansion makes that class of error structurally impossible rather than something
to be prompted against.

\subsection{Shared guardrails}

The same failure threatens both query languages, so the same two transforms are
applied to both, identically. Any trailing \texttt{LIMIT} in a generated query is
raised to 2{,}000 rows before execution; the result set is then balanced to at
most 6 rows per entity and 48 rows overall, where the entity column is detected by
name. Applying these to SQL and Cypher alike means the comparison measures
retrieval architecture rather than which prompt happened to dodge the trap.

Both query-writing arms also get exactly one repair attempt, triggered when a
query raises an exception or returns zero rows: the model sees the error text and
rewrites once. There is no second repair. \Cref{sec:results} shows the two arms
used this allowance very unequally.

\subsection{Model access}

All three arms call a single OpenAI-compatible endpoint provided by Bifrost, an
open-source LLM gateway, which routes to DeepSeek or to AWS Bedrock. The
benchmarked configuration is \texttt{deepseek-v4-flash} at temperature 0 with
reasoning disabled, 700 output tokens for query generation and 1{,}500 for
answers.

Routing through one endpoint is a methodological device, not a convenience. It
makes ``all three arms used the same model with the same settings'' a structural
property of the deployment rather than a convention maintained by discipline: the
model is a request parameter read from one place by all three arms. It also makes
the cross-model probe in \Cref{sec:failures} a configuration change rather than a
code change. Reasoning was disabled because an uncapped reasoning trace consumed
roughly 33{,}000 completion tokens per call, which \texttt{max\_tokens} does not
bound; the setting is identical for every arm, so the comparison is unaffected.
